\documentclass[../../main.tex]{subfiles}

\begin{document}

\section{度量拓扑}\label{section:ch02-度量拓扑}

在集合上定义拓扑，最重要最常用的方法之一就是藉助于这个集合的度量来实现。用这种方式给出拓扑是现代分析的核心之一。这一节我们将定义度量拓扑，并且将给出大量例子。下一节再研究度量拓扑所满足的一些性质。

\begin{definition}[度量与 $\varepsilon$-球]\label{definition:ch02-度量}
集合 $X$ 的一个度量（metric）是一个函数
\begin{align*}
d: X \times X \longrightarrow \mathbb{R}
\end{align*}
使得以下性质成立：
\begin{enumerate}[(1)]
\item 对于所有的 $x, y \in X$，$d(x, y) \geqslant 0$；等号当且仅当 $x = y$ 时成立。
\item 对于所有 $x, y \in X$，$d(x, y) = d(y, x)$。
\item （三角不等式）对于所有 $x, y, z \in X$，$d(x, y) + d(y, z) \geqslant d(x, z)$。
\end{enumerate}

给定 $X$ 的一个度量 $d$，数 $d(x, y)$ 通常称为 $x$ 与 $y$ 之间在度量 $d$ 下的距离（distance）。对于 $\varepsilon > 0$，考虑所有与 $x$ 的距离小于 $\varepsilon$ 的点 $y$ 的集合
\begin{align*}
B_{d}(x, \varepsilon) = \{y \mid d(x, y) < \varepsilon\},
\end{align*}
它称为以 $x$ 为中心的 $\varepsilon$-球（$\varepsilon$-ball centered at $x$）。在不致引起混淆的情况下，有时在记号中省略度量 $d$，而把这个球简单地记为 $B(x, \varepsilon)$。

若 $d$ 是集合 $X$ 的一个度量，则全体 $\varepsilon$-球 $B(x, \varepsilon)$ 的族，其中 $x \in X$，$\varepsilon > 0$，是 $X$ 的某一个拓扑的基，这个拓扑称为由度量 $d$ 诱导出来的度量拓扑（metric topology）。
\end{definition}

因为对于任意 $\varepsilon > 0$ 有 $x \in B(x, \varepsilon)$，所以基定义中的第一个条件显然满足。在验证基定义中的第二个条件之前，我们首先证明：若 $y$ 是基元素 $B(x, \varepsilon)$ 的一个点，则存在以 $y$ 为中心的一个基元素 $B(y, \delta)$，它包含在 $B(x, \varepsilon)$ 之中。取 $\delta$ 为正数 $\varepsilon - d(x, y)$，则当 $z \in B(y, \delta)$ 时，$d(y, z) < \varepsilon - d(x, y)$，由此得到
\begin{align*}
d(x, z) \leqslant d(x, y) + d(y, z) < \varepsilon,
\end{align*}
从而证明了 $B(y, \delta) \subset B(x, \varepsilon)$。参见\cref{figure:ch02-度量拓扑-20.1}。

现在来验证基的第二个条件。设 $B_{1}$ 和 $B_{2}$ 为两个基元素，并且 $y \in B_{1} \cap B_{2}$。根据刚才的证明，可以取正数 $\delta_{1}$ 和 $\delta_{2}$，使得 $B(y, \delta_{1}) \subset B_{1}$ 和 $B(y, \delta_{2}) \subset B_{2}$。令 $\delta$ 为 $\delta_{1}$ 和 $\delta_{2}$ 中较小的一个，便可见 $B(y, \delta) \subset B_{1} \cap B_{2}$。

我们应用刚才证明的结论，将度量拓扑的定义重新改写成：集合 $U$ 是由 $d$ 诱导出来的度量拓扑中的开集，当且仅当对于每一个 $y \in U$，存在一个 $\delta > 0$ 使得 $B_{d}(y, \delta) \subset U$。

这个条件显然蕴涵 $U$ 是开集，反之，若 $U$ 是开集，它必定包含着包含点 $y$ 的一个基元素 $B = B_{d}(x, \varepsilon)$。因此，$B$ 又包含着以 $y$ 为中心的一个基元素 $B_{d}(y, \delta)$。

\par
\begin{figure}[htbp]
\centering
\includegraphics[width=0.4\textwidth]{ch02-fig20-1}
\caption{$\varepsilon$-球与基条件}\label{figure:ch02-度量拓扑-20.1}
\end{figure}

\begin{examples}\label{examples:ch02-度量拓扑-例}
\begin{example}\label{example:ch02-离散度量}
对于集合 $X$，定义
\begin{align*}
d(x, y) = 1 \quad \text{如果 } x \neq y, \qquad d(x, y) = 0 \quad \text{如果 } x = y.
\end{align*}
容易验证 $d$ 是一个度量。由 $d$ 诱导出来的拓扑是离散拓扑。例如，基元素 $B(x, 1)$ 就只包含着点 $x$。
\end{example}

\begin{example}\label{example:ch02-标准度量}
用等式
\begin{align*}
d(x, y) = |x - y|
\end{align*}
所定义的度量是实直线 $\mathbb{R}$ 上的标准度量。容易验证 $d$ 是一个度量。由它诱导出来的拓扑与序拓扑相同。这是因为序拓扑中的每一个基元素都是度量拓扑中的一个基元素。事实上，
\begin{align*}
(a, b) = B(x, \varepsilon),
\end{align*}
其中 $x = (a + b)/2$，$\varepsilon = (b - a)/2$。反之，每一个 $\varepsilon$-球 $B(x, \varepsilon)$ 等于一个开区间 $(x - \varepsilon, x + \varepsilon)$。
\end{example}
\end{examples}

\begin{definition}[可度量化空间与有界集]\label{definition:ch02-可度量化有界集}
设 $X$ 是一个拓扑空间。如果 $X$ 的拓扑是由集合 $X$ 的某一个度量 $d$ 所诱导出来的，则称 $X$ 是一个可度量化（metrizable）空间。度量空间（metric space）指的就是一个可度量化空间 $X$，连同一个诱导出 $X$ 的拓扑的特定的度量 $d$。

数学中许多重要的空间都是可度量化的，但也有一些不是。对于一个空间来说，可度量化往往是一个最理想的性质，因为度量的存在性为空间中一些定理的证明提供了一个重要工具。

因此，拓扑学中一个基本的也是重要的课题就是寻找保证拓扑空间可度量化的条件。我们在第 4 章中的目标之一，就是求出这样的条件，它便是著名的 Urysohn 度量化定理。更进一步的度量化定理在第 6 章讲述。本节只限于证明 $\mathbb{R}^{n}$ 和 $\mathbb{R}^{\omega}$ 是可度量化的。

虽然可度量化是拓扑学的一个重要课题，但是对于度量空间的研究，并不真正属于拓扑学，而更多的是属于分析学。一个空间是否是可度量化的，仅仅依赖于空间的拓扑。但与 $X$ 的具体度量有关的一些性质一般并不是这样，所以它们不是拓扑性质。例如，我们可以在度量空间中给出下面的定义。

设 $X$ 是具有度量 $d$ 的一个度量空间。$X$ 的子集 $A$ 称为有界的，如果存在某数 $M$，使得对于 $A$ 中任意两点 $a_{1}$ 和 $a_{2}$ 有
\begin{align*}
d(a_{1}, a_{2}) \leqslant M.
\end{align*}
若 $A$ 是有界的，$A$ 的直径（diameter）定义为一个数
\begin{align*}
\operatorname{diam} A = \sup \{d(a_{1}, a_{2}) \mid a_{1}, a_{2} \in A\}.
\end{align*}
\end{definition}

集合的有界性就不是一个拓扑性质，这是由于它仅依赖于 $X$ 所采用的特定度量。例如，若 $X$ 是具有度量 $d$ 的一个度量空间，则存在一个度量 $\bar{d}$ 诱导出的 $X$ 的拓扑，使得对于 $\bar{d}$ 而言，$X$ 的每一个子集都是有界的。$\bar{d}$ 可以由以下定理给出。

\begin{theorem}[标准有界度量]\label{theorem:ch02-标准有界度量}
设 $X$ 是具有度量 $d$ 的一个度量空间。则用
\begin{align*}
\bar{d}(x, y) = \min \{d(x, y), 1\}
\end{align*}
所定义的 $\bar{d}: X \times X \rightarrow \mathbb{R}$ 是一个度量，并且 $\bar{d}$ 和 $d$ 诱导出的拓扑是 $X$ 的同一个拓扑。

度量 $\bar{d}$ 称为相应于 $d$ 的标准有界度量（standard bounded metric）。
\end{theorem}

\begin{proof}
容易验证 $\bar{d}$ 满足度量的前两个条件。下面验证三角不等式，即
\begin{align*}
\bar{d}(x, z) \leqslant \bar{d}(x, y) + \bar{d}(y, z).
\end{align*}
若 $d(x, y) \geqslant 1$ 或 $d(y, z) \geqslant 1$，则不等式右端至少为 1；而根据定义，其左边最多为 1，所以不等式成立。下面只考虑 $d(x, y) < 1$ 且 $d(y, z) < 1$ 的情形。这时有
\begin{align*}
d(x, z) \leqslant d(x, y) + d(y, z) = \bar{d}(x, y) + \bar{d}(y, z).
\end{align*}
根据定义，$\bar{d}(x, z) \leqslant d(x, z)$，所以对于 $\bar{d}$，三角不等式成立。

注意在任何一个度量空间中，满足 $\varepsilon < 1$ 的 $\varepsilon$-球形成度量拓扑的一个基，这是由于含有 $x$ 的每一基元素都包含一个这种在点 $x$ 的 $\varepsilon$-球。因为对于两种度量 $d$ 和 $\bar{d}$ 而言，满足 $\varepsilon < 1$ 的 $\varepsilon$-球的集合是同一个集合，从而这两种度量诱导出的是 $X$ 上的同一个拓扑。
\end{proof}

现在来考虑一些我们所熟悉的空间，并证明它们是可度量化的。

\begin{definition}[欧氏度量与平方度量]\label{definition:ch02-欧氏平方度量}
给定 $\mathbb{R}^{n}$ 中的点 $\boldsymbol{x} = (x_{1}, \cdots, x_{n})$。$\boldsymbol{x}$ 的模（norm）定义为
\begin{align*}
\|\boldsymbol{x}\| = (x_{1}^{2} + \cdots + x_{n}^{2})^{1/2}.
\end{align*}
$\mathbb{R}^{n}$ 中的欧氏度量（Euclidean metric）$d$ 定义为
\begin{align*}
d(\boldsymbol{x}, \boldsymbol{y}) = \|\boldsymbol{x} - \boldsymbol{y}\| = [(x_{1} - y_{1})^{2} + \cdots + (x_{n} - y_{n})^{2}]^{1/2}.
\end{align*}
平方度量（square metric）$\rho$ 定义为
\begin{align*}
\rho(\boldsymbol{x}, \boldsymbol{y}) = \max \{|x_{1} - y_{1}|, \cdots, |x_{n} - y_{n}|\}.
\end{align*}
\end{definition}

证明 $d$ 是一个度量还要做一些工作。如果读者还不清楚这一点的话，可以在习题中找到证明的要点。我们很少有机会用到 $\mathbb{R}^{n}$ 的欧氏度量。

证明 $\rho$ 是一个度量是很容易的，只有三角不等式不那么明显。由 $\mathbb{R}$ 中的三角不等式得知，对于每一个正整数 $i$，有
\begin{align*}
|x_{i} - z_{i}| \leqslant |x_{i} - y_{i}| + |y_{i} - z_{i}|.
\end{align*}
于是根据 $\rho$ 的定义有
\begin{align*}
|x_{i} - z_{i}| \leqslant \rho(\boldsymbol{x}, \boldsymbol{y}) + \rho(\boldsymbol{y}, \boldsymbol{z}).
\end{align*}
最后得到
\begin{align*}
\rho(\boldsymbol{x}, \boldsymbol{z}) = \max \{|x_{i} - z_{i}|\} \leqslant \rho(\boldsymbol{x}, \boldsymbol{y}) + \rho(\boldsymbol{y}, \boldsymbol{z}).
\end{align*}

在实直线 $\mathbb{R} = \mathbb{R}^{1}$ 上，这两个度量与 $\mathbb{R}$ 的标准度量相同。对于平面 $\mathbb{R}^{2}$，关于度量 $d$ 的基元素可以画成圆域，关于度量 $\rho$ 的基元素可以画成方域。

我们将要证明这两种度量都诱导出 $\mathbb{R}^{n}$ 上的通常拓扑。为此需要以下引理。

\begin{lemma}[两度量间的粗细比较]\label{lemma:ch02-度量粗细比较}
设 $d$ 和 $d'$ 是集合 $X$ 上的两个度量。$\mathcal{T}$ 和 $\mathcal{T}'$ 分别是由它们诱导的拓扑，则 $\mathcal{T}'$ 细于 $\mathcal{T}$ 当且仅当对于每一个 $x \in X$ 及每一个 $\varepsilon > 0$，存在一个 $\delta > 0$，使得
\begin{align*}
B_{d'}(x, \delta) \subset B_{d}(x, \varepsilon).
\end{align*}
\end{lemma}

\begin{proof}
设 $\mathcal{T}'$ 细于 $\mathcal{T}$。给定 $\mathcal{T}$ 的一个基元素 $B_{d}(x, \varepsilon)$，根据\cref{lemma:ch02-开集族为基}，存在拓扑 $\mathcal{T}'$ 的一个基元素 $B'$，使得 $x \in B' \subset B_{d}(x, \varepsilon)$，于是存在一个以 $x$ 为中心的球 $B_{d'}(x, \delta) \subset B'$。

反之，假定 $\varepsilon$-$\delta$ 条件成立。对于 $\mathcal{T}$ 中含有点 $x$ 的一个基元素 $B$，可以求出 $B$ 中以 $x$ 为中心的一个球 $B_{d}(x, \varepsilon)$。再根据已知条件，存在一个 $\delta$，使得 $B_{d'}(x, \delta) \subset B_{d}(x, \varepsilon)$。于是根据\cref{lemma:ch02-开集族为基}得到 $\mathcal{T}'$ 细于 $\mathcal{T}$。
\end{proof}

\begin{theorem}[欧氏度量、平方度量与积拓扑一致]\label{theorem:ch02-Rn度量拓扑}
由欧氏度量 $d$ 及平方度量 $\rho$ 所诱导的 $\mathbb{R}^{n}$ 的拓扑与 $\mathbb{R}^{n}$ 的积拓扑相同。
\end{theorem}

\begin{proof}
设 $\boldsymbol{x} = (x_{1}, \cdots, x_{n})$ 与 $\boldsymbol{y} = (y_{1}, \cdots, y_{n})$ 是 $\mathbb{R}^{n}$ 的两个点。可以通过简单的代数运算验证
\begin{align*}
\rho(\boldsymbol{x}, \boldsymbol{y}) \leqslant d(\boldsymbol{x}, \boldsymbol{y}) \leqslant \sqrt{n}\, \rho(\boldsymbol{x}, \boldsymbol{y}).
\end{align*}
上式中前一个不等式表明，
\begin{align*}
B_{d}(\boldsymbol{x}, \varepsilon) \subset B_{\rho}(\boldsymbol{x}, \varepsilon)
\end{align*}
对于所有的 $\boldsymbol{x}$ 和 $\varepsilon$ 成立，这是因为若 $d(\boldsymbol{x}, \boldsymbol{y}) < \varepsilon$，则 $\rho(\boldsymbol{x}, \boldsymbol{y}) < \varepsilon$。类似地，后一个不等式表明，对于所有的 $\boldsymbol{x}$ 及 $\varepsilon$ 有
\begin{align*}
B_{\rho}(\boldsymbol{x}, \varepsilon/\sqrt{n}) \subset B_{d}(\boldsymbol{x}, \varepsilon).
\end{align*}
根据上一个引理可见，这两个度量拓扑相同。

现在来证明积拓扑与由度量 $\rho$ 所给出的拓扑是相同的。首先，令
\begin{align*}
B = (a_{1}, b_{1}) \times \cdots \times (a_{n}, b_{n})
\end{align*}
为积拓扑中的一个基元素，$\boldsymbol{x} = (x_{1}, \cdots, x_{n})$ 为 $B$ 的一个元素。对于每一个 $i$，存在一个 $\varepsilon_{i}$，使得
\begin{align*}
(x_{i} - \varepsilon_{i}, x_{i} + \varepsilon_{i}) \subset (a_{i}, b_{i});
\end{align*}
选取 $\varepsilon = \min \{\varepsilon_{1}, \cdots, \varepsilon_{n}\}$。容易验证 $B_{\rho}(\boldsymbol{x}, \varepsilon) \subset B$。这就证明了 $\rho$-拓扑细于积拓扑。

反之，设 $B_{\rho}(\boldsymbol{x}, \varepsilon)$ 是 $\rho$-拓扑的一个基元素，给定元素 $\boldsymbol{y} \in B_{\rho}(\boldsymbol{x}, \varepsilon)$，我们需要找出积拓扑的一个基元素 $B$，使得
\begin{align*}
\boldsymbol{y} \in B \subset B_{\rho}(\boldsymbol{x}, \varepsilon).
\end{align*}
然而，这是显然的。因为，对于积拓扑而言，
\begin{align*}
B_{\rho}(\boldsymbol{x}, \varepsilon) = (x_{1} - \varepsilon, x_{1} + \varepsilon) \times \cdots \times (x_{n} - \varepsilon, x_{n} + \varepsilon)
\end{align*}
本身就是一个基元素。
\end{proof}

现在我们来考虑无限笛卡儿积 $\mathbb{R}^{\omega}$。人们自然会想到将度量 $d$ 及度量 $\rho$ 推广到这个空间。比如，我们尝试用
\begin{align*}
d(\boldsymbol{x}, \boldsymbol{y}) = \left[\sum_{i=1}^{\infty} (x_{i} - y_{i})^{2}\right]^{1/2}
\end{align*}
定义 $\mathbb{R}^{\omega}$ 的一个度量 $d$。但是这一公式未必有意义，原因是级数不一定收敛。（然而，这个公式可以在 $\mathbb{R}^{\omega}$ 的某一个重要子集上定义一个度量，见习题。）

类似地，人们也试图用
\begin{align*}
\rho(\boldsymbol{x}, \boldsymbol{y}) = \sup \{|x_{n} - y_{n}|\}
\end{align*}
将平方度量 $\rho$ 推广到 $\mathbb{R}^{\omega}$ 上。但是，这再度涉及到它未必有意义的问题。然而，若我们用 $\mathbb{R}$ 中相应的有界度量 $\bar{d}(x, y) = \min \{|x - y|, 1\}$ 来代替度量 $d(x, y) = |x - y|$，则定义有意义。由此引出的 $\mathbb{R}^{\omega}$ 上的度量称为一致度量。

对于任意指标集 $J$，可按以下方式在一般的笛卡儿积 $\mathbb{R}^{J}$ 上定义一致度量如下：

\begin{definition}[一致度量与一致拓扑]\label{definition:ch02-一致度量}
给定指标集 $J$ 以及 $\mathbb{R}^{J}$ 中的点 $\boldsymbol{x} = (x_{\alpha})_{\alpha \in J}$ 和 $\boldsymbol{y} = (y_{\alpha})_{\alpha \in J}$，定义 $\mathbb{R}^{J}$ 的一个度量 $\bar{\rho}$ 为
\begin{align*}
\bar{\rho}(\boldsymbol{x}, \boldsymbol{y}) = \sup \{\bar{d}(x_{\alpha}, y_{\alpha}) \mid \alpha \in J\},
\end{align*}
其中 $\bar{d}$ 是 $\mathbb{R}$ 的标准有界度量，$\bar{\rho}$ 称为 $\mathbb{R}^{J}$ 上的一致度量（uniform metric），由 $\bar{\rho}$ 所诱导出来的拓扑称为一致拓扑（uniform topology）。
\end{definition}

一致拓扑与积拓扑和箱拓扑之间有以下关系：

\begin{theorem}[一致拓扑的位置]\label{theorem:ch02-一致拓扑位置}
$\mathbb{R}^{J}$ 上的一致拓扑细于积拓扑，粗于箱拓扑。当 $J$ 为无限集时，这三个拓扑两两不同。
\end{theorem}

\begin{proof}
设给定一个点 $\boldsymbol{x} = (x_{\alpha})_{\alpha \in J}$ 和包含 $\boldsymbol{x}$ 的一个积拓扑基元素 $\prod U_{\alpha}$。令 $\alpha_{1}, \cdots, \alpha_{n}$ 是使 $U_{\alpha} \neq \mathbb{R}$ 的那些指标。对于每一个 $\alpha_{i}$，选取 $\varepsilon_{i} > 0$，使得以 $x_{\alpha_{i}}$ 为中心关于度量 $\bar{d}$ 的 $\varepsilon_{i}$-球包含于 $U_{\alpha_{i}}$。由于 $U_{\alpha_{i}}$ 在 $\mathbb{R}$ 中是开的，所以这样的 $\varepsilon_{i}$-球是存在的。令 $\varepsilon = \min \{\varepsilon_{1}, \cdots, \varepsilon_{n}\}$。则有以 $x_{\alpha_{i}}$ 为中心关于度量 $\bar{\rho}$ 的 $\varepsilon$-球包含于 $\prod U_{\alpha}$。对于 $\mathbb{R}^{J}$ 的一个点 $\boldsymbol{z}$，若 $\bar{\rho}(\boldsymbol{x}, \boldsymbol{z}) < \varepsilon$，则对于所有 $\alpha$，$\bar{d}(x_{\alpha}, z_{\alpha}) < \varepsilon$，因此 $\boldsymbol{z} \in \prod U_{\alpha}$。从而，一致拓扑细于积拓扑。

另一方面，令 $B$ 为以 $\boldsymbol{x}$ 为中心关于度量 $\bar{\rho}$ 的 $\varepsilon$-球。这时，$\boldsymbol{x}$ 在箱拓扑下的邻域
\begin{align*}
U = \prod_{\alpha} \left(x_{\alpha} - \frac{1}{2}\varepsilon, x_{\alpha} + \frac{1}{2}\varepsilon\right)
\end{align*}
包含于 $B$。这是因为，若 $\boldsymbol{y} \in U$，则对于所有的 $\alpha$ 有 $\bar{d}(x_{\alpha}, y_{\alpha}) < \frac{1}{2}\varepsilon$，因此 $\bar{\rho}(\boldsymbol{x}, \boldsymbol{y}) \leqslant \frac{1}{2}\varepsilon$。

当 $J$ 是无限集时，这三个拓扑互不相同的证明我们留作习题。
\end{proof}

对于 $J$ 为无限集的情形，我们还没能决定 $\mathbb{R}^{J}$ 上的箱拓扑与积拓扑是否是可度量化的。事实上，我们将看到，只是在 $J$ 是可数集并且取积拓扑的情况下，$\mathbb{R}^{J}$ 才是可度量化的。

\begin{theorem}[诱导 $\mathbb{R}^{\omega}$ 积拓扑的度量]\label{theorem:ch02-Romega度量}
设 $\bar{d}(a, b) = \min \{|a - b|, 1\}$ 是 $\mathbb{R}$ 上的标准有界度量。对于 $\mathbb{R}^{\omega}$ 的两个点 $\boldsymbol{x}, \boldsymbol{y}$，定义
\begin{align*}
D(\boldsymbol{x}, \boldsymbol{y}) = \sup \left\{\frac{\bar{d}(x_{i}, y_{i})}{i}\right\}.
\end{align*}
那么 $D$ 是诱导 $\mathbb{R}^{\omega}$ 的积拓扑的一个度量。
\end{theorem}

\begin{proof}
除了三角不等式外，度量的其他性质显然是满足的。以下证明三角不等式，注意对于所有的 $i$ 有
\begin{align*}
\frac{\bar{d}(x_{i}, z_{i})}{i} \leqslant \frac{\bar{d}(x_{i}, y_{i})}{i} + \frac{\bar{d}(y_{i}, z_{i})}{i} \leqslant D(\boldsymbol{x}, \boldsymbol{y}) + D(\boldsymbol{y}, \boldsymbol{z}),
\end{align*}
因而
\begin{align*}
\sup \left\{\frac{\bar{d}(x_{i}, z_{i})}{i}\right\} \leqslant D(\boldsymbol{x}, \boldsymbol{y}) + D(\boldsymbol{y}, \boldsymbol{z}).
\end{align*}

要证明 $D$ 诱导出积拓扑还得花点力气。首先，令 $U$ 为度量拓扑中的一个开集，$\boldsymbol{x} \in U$。在积拓扑中找一个开集 $V$，使得 $\boldsymbol{x} \in V \subset U$。在 $U$ 中选取一个 $\varepsilon$-球 $B_{D}(\boldsymbol{x}, \varepsilon)$。其次，可将 $N$ 取得足够大，使得 $\frac{1}{N} < \varepsilon$。最后，令 $V$ 是积拓扑的基元素
\begin{align*}
V = (x_{1} - \varepsilon, x_{1} + \varepsilon) \times \cdots \times (x_{N} - \varepsilon, x_{N} + \varepsilon) \times \mathbb{R} \times \mathbb{R} \times \cdots.
\end{align*}
我们来证明 $V \subset B_{D}(\boldsymbol{x}, \varepsilon)$：对于 $\mathbb{R}^{\omega}$ 中的任意一点 $\boldsymbol{y}$ 以及 $i \geqslant N$，
\begin{align*}
\frac{\bar{d}(x_{i}, y_{i})}{i} \leqslant \frac{1}{N}.
\end{align*}
因此
\begin{align*}
D(\boldsymbol{x}, \boldsymbol{y}) \leqslant \max \left\{\frac{\bar{d}(x_{1}, y_{1})}{1}, \cdots, \frac{\bar{d}(x_{N}, y_{N})}{N}, \frac{1}{N}\right\}.
\end{align*}
如果 $\boldsymbol{y}$ 在 $V$ 中，上面的式子小于 $\varepsilon$，于是 $V \subset B_{D}(\boldsymbol{x}, \varepsilon)$，这便是所要证的。

反之，考虑积拓扑的一个基元素
\begin{align*}
U = \prod_{i \in \mathbb{Z}_{+}} U_{i},
\end{align*}
其中，当 $i = \alpha_{1}, \cdots, \alpha_{n}$ 时，$U_{i}$ 是 $\mathbb{R}$ 中的开集。对于其他的指标 $i$，$U_{i} = \mathbb{R}$。给定 $\boldsymbol{x} \in U$，可以在度量拓扑中选取一个开集 $V$，使得 $\boldsymbol{x} \in V \subset U$。在 $\mathbb{R}$ 中选取一个以 $x_{i}$ 为中心并且被 $U_{i}$ 所包含的区间 $(x_{i} - \varepsilon_{i}, x_{i} + \varepsilon_{i})$，其中，$i = \alpha_{1}, \cdots, \alpha_{n}$，要求 $\varepsilon_{i} \leqslant 1$。定义
\begin{align*}
\varepsilon = \min \{\varepsilon_{i}/i \mid i = \alpha_{1}, \cdots, \alpha_{n}\}.
\end{align*}
我们来证明
\begin{align*}
\boldsymbol{x} \in B_{D}(\boldsymbol{x}, \varepsilon) \subset U.
\end{align*}
事实上，设 $\boldsymbol{y}$ 是 $B_{D}(\boldsymbol{x}, \varepsilon)$ 的一个点，那么对于所有的 $i$，
\begin{align*}
\frac{\bar{d}(x_{i}, y_{i})}{i} \leqslant D(\boldsymbol{x}, \boldsymbol{y}) < \varepsilon.
\end{align*}
若 $i = \alpha_{1}, \cdots, \alpha_{n}$，则 $\varepsilon \leqslant \varepsilon_{i}/i$，因此 $\bar{d}(x_{i}, y_{i}) < \varepsilon_{i} \leqslant 1$。由此推出 $|x_{i} - y_{i}| < \varepsilon_{i}$。因此，$\boldsymbol{y} \in \prod U_{i}$。
\end{proof}

\subsection*{习题}

\begin{exercise}\label{exercise:ch02-度量拓扑-1}
(a) 在 $\mathbb{R}^{n}$ 中定义
\begin{align*}
d'(\boldsymbol{x}, \boldsymbol{y}) = |x_{1} - y_{1}| + \cdots + |x_{n} - y_{n}|.
\end{align*}
证明 $d'$ 是诱导出 $\mathbb{R}^{n}$ 的通常拓扑的一个度量。当 $n = 2$ 时，画出 $d'$ 下的基元素。

(b) 更一般地，对于 $p \geqslant 1$ 和 $\boldsymbol{x}, \boldsymbol{y} \in \mathbb{R}^{n}$，定义
\begin{align*}
d'(\boldsymbol{x}, \boldsymbol{y}) = \left[\sum_{i=1}^{n} |x_{i} - y_{i}|^{p}\right]^{1/p},
\end{align*}
假定 $d'$ 是一个度量。试证明它诱导出 $\mathbb{R}^{n}$ 的通常拓扑。
\end{exercise}

\begin{exercise}\label{exercise:ch02-度量拓扑-2}
证明：$\mathbb{R} \times \mathbb{R}$ 在字典序拓扑下是可度量化的。
\end{exercise}

\begin{exercise}\label{exercise:ch02-度量拓扑-3}
设 $X$ 是以 $d$ 为度量的一个度量空间。

(a) 证明 $d: X \times X \rightarrow \mathbb{R}$ 是连续的。

(b) 设 $X'$ 是一个空间，作为集合与 $X$ 相同。证明：若 $d: X' \times X' \rightarrow \mathbb{R}$ 连续，则 $X'$ 的拓扑细于 $X$ 的拓扑。

可将这个习题中的结论总结为：若 $X$ 有一度量 $d$，则 $d$ 诱导的拓扑是使得函数 $d$ 连续的所有拓扑中最粗的拓扑。
\end{exercise}

\begin{exercise}\label{exercise:ch02-度量拓扑-4}
考虑 $\mathbb{R}^{\omega}$ 的积拓扑、一致拓扑和箱拓扑。

(a) 哪些拓扑是使得以下从 $\mathbb{R}$ 到 $\mathbb{R}^{\omega}$ 的函数连续的拓扑？
\begin{align*}
f(t) = (t, 2t, 3t, \cdots), \qquad g(t) = (t, t, t, \cdots), \qquad h(t) = \left(t, \frac{1}{2}t, \frac{1}{3}t, \cdots\right).
\end{align*}

(b) 哪些拓扑是使得以下序列收敛的拓扑？
\begin{align*}
\boldsymbol{w}_{1} = (1, 1, 1, 1, \cdots), \quad \boldsymbol{x}_{1} = (1, 1, 1, 1, \cdots), \\
\boldsymbol{w}_{2} = (0, 2, 2, 2, \cdots), \quad \boldsymbol{x}_{2} = \left(0, \frac{1}{2}, \frac{1}{2}, \frac{1}{2}, \cdots\right), \\
\boldsymbol{w}_{3} = (0, 0, 3, 3, \cdots), \quad \boldsymbol{x}_{3} = \left(0, 0, \frac{1}{3}, \frac{1}{3}, \cdots\right), \\
\cdots \quad \cdots \\
\boldsymbol{y}_{1} = (1, 0, 0, 0, \cdots), \quad \boldsymbol{z}_{1} = (1, 1, 0, 0, \cdots), \\
\boldsymbol{y}_{2} = \left(\frac{1}{2}, \frac{1}{2}, 0, 0, \cdots\right), \quad \boldsymbol{z}_{2} = \left(\frac{1}{2}, \frac{1}{2}, 0, 0, \cdots\right), \\
\boldsymbol{y}_{3} = \left(\frac{1}{3}, \frac{1}{3}, \frac{1}{3}, 0, \cdots\right), \quad \boldsymbol{z}_{3} = \left(\frac{1}{3}, \frac{1}{3}, 0, 0, \cdots\right), \\
\cdots \quad \cdots
\end{align*}
\end{exercise}

\begin{exercise}\label{exercise:ch02-度量拓扑-5}
设 $\mathbb{R}^{\infty}$ 是由 $\mathbb{R}^{\omega}$ 中所有终端为 0 的序列组成的子集。在一致拓扑下，$\mathbb{R}^{\infty}$ 在 $\mathbb{R}^{\omega}$ 中的闭包是什么？证明你的结论。
\end{exercise}

\begin{exercise}\label{exercise:ch02-度量拓扑-6}
设 $\bar{\rho}$ 是 $\mathbb{R}^{\omega}$ 的一致度量。给定 $\boldsymbol{x} = (x_{1}, x_{2}, \cdots) \in \mathbb{R}^{\omega}$ 和 $0 < \varepsilon < 1$，令
\begin{align*}
U(\boldsymbol{x}, \varepsilon) = (x_{1} - \varepsilon, x_{1} + \varepsilon) \times \cdots \times (x_{n} - \varepsilon, x_{n} + \varepsilon) \times \cdots.
\end{align*}
(a) 证明 $U(\boldsymbol{x}, \varepsilon)$ 不等于 $\varepsilon$-球 $B_{\bar{\rho}}(\boldsymbol{x}, \varepsilon)$。

(b) 证明 $U(\boldsymbol{x}, \varepsilon)$ 不是一致拓扑空间中的开集。

(c) 证明
\begin{align*}
B_{\bar{\rho}}(\boldsymbol{x}, \varepsilon) = \bigcup_{\delta < \varepsilon} U(\boldsymbol{x}, \delta).
\end{align*}
\end{exercise}

\begin{exercise}\label{exercise:ch02-度量拓扑-7}
考虑第 19 节习题 8 中定义的映射 $h: \mathbb{R}^{\omega} \rightarrow \mathbb{R}^{\omega}$，并且赋予 $\mathbb{R}^{\omega}$ 一致拓扑。当 $a_{i}$ 和 $b_{i}$ 满足什么条件时，$h$ 连续？什么条件会使得 $h$ 为同胚？
\end{exercise}

\begin{exercise}\label{exercise:ch02-度量拓扑-8}
设 $X$ 为 $\mathbb{R}^{\omega}$ 中所有满足 $\sum x_{i}^{2}$ 收敛的序列 $\boldsymbol{x}$ 所构成的子集。用公式
\begin{align*}
d(\boldsymbol{x}, \boldsymbol{y}) = \left[\sum_{i=1}^{\infty} (x_{i} - y_{i})^{2}\right]^{1/2}
\end{align*}
定义 $X$ 上的一个度量（见习题 10）。$X$ 分别由 $\mathbb{R}^{\omega}$ 的箱拓扑、一致拓扑和积拓扑继承了三个拓扑。用上式给出的度量 $d$ 也定义了 $X$ 的一个拓扑，称之为 $\ell^{2}$-拓扑。

(a) 证明：对于 $X$ 而言，以下包含关系成立：
\begin{align*}
\text{箱拓扑} \supset \ell^{2}\text{-拓扑} \supset \text{一致拓扑}.
\end{align*}

(b) 所有终端为 0 的序列所构成的集合 $\mathbb{R}^{\infty}$ 是 $X$ 的子集。证明：作为 $X$ 的子空间，$\mathbb{R}^{\infty}$ 所继承的四种拓扑互不相同。

(c) 集合
\begin{align*}
H = \prod_{n \in \mathbb{Z}_{+}} [0, 1/n]
\end{align*}
是 $X$ 的一个子集。称之为 Hilbert 立方（Hilbert cube）。试比较 $H$ 作为 $X$ 的子空间所继承的四种拓扑。
\end{exercise}

\begin{exercise}\label{exercise:ch02-度量拓扑-9}
我们可以这样证明 $\mathbb{R}^{n}$ 的欧氏度量 $d$ 是一个度量：对于 $\boldsymbol{x}, \boldsymbol{y} \in \mathbb{R}^{n}$ 和 $c \in \mathbb{R}$，定义
\begin{align*}
\boldsymbol{x} + \boldsymbol{y} = (x_{1} + y_{1}, \cdots, x_{n} + y_{n}), \\
c\boldsymbol{x} = (cx_{1}, \cdots, cx_{n}), \\
\boldsymbol{x} \cdot \boldsymbol{y} = x_{1}y_{1} + \cdots + x_{n}y_{n}.
\end{align*}
(a) 证明 $\boldsymbol{x} \cdot (\boldsymbol{y} + \boldsymbol{z}) = (\boldsymbol{x} \cdot \boldsymbol{y}) + (\boldsymbol{x} \cdot \boldsymbol{z})$。

(b) 证明 $|\boldsymbol{x} \cdot \boldsymbol{y}| \leqslant \|\boldsymbol{x}\|\, \|\boldsymbol{y}\|$。[提示：若 $\boldsymbol{x}, \boldsymbol{y} \neq 0$，令 $a = 1/\|\boldsymbol{x}\|$，$b = 1/\|\boldsymbol{y}\|$，应用 $\|a\boldsymbol{x} \pm b\boldsymbol{y}\| \geqslant 0$。]

(c) 证明 $\|\boldsymbol{x} + \boldsymbol{y}\| \leqslant \|\boldsymbol{x}\| + \|\boldsymbol{y}\|$。[提示：计算 $(\boldsymbol{x} + \boldsymbol{y}) \cdot (\boldsymbol{x} + \boldsymbol{y})$，并应用 (b) 小题。]

(d) 证明 $d$ 是一个度量。
\end{exercise}

\begin{exercise}\label{exercise:ch02-度量拓扑-10}
设 $X$ 为 $\mathbb{R}^{\omega}$ 中所有使得 $\sum x_{i}^{2}$ 收敛的序列 $(x_{1}, x_{2}, \cdots)$ 所构成的集合。（这里我们要承认有关无穷数列的某些基本性质。倘若你对这些性质不熟悉，请参见下一节的习题 11。）

(a) 证明：若 $\boldsymbol{x}, \boldsymbol{y} \in X$，则 $\sum |x_{i}y_{i}|$ 收敛。[提示：应用习题 9 中 (b) 小题，证明部分和有界。]

(b) 设 $c \in \mathbb{R}$。证明：若 $\boldsymbol{x}, \boldsymbol{y} \in X$，则有 $\boldsymbol{x} + \boldsymbol{y} \in X$ 和 $c\boldsymbol{x} \in X$。

(c) 证明
\begin{align*}
d(\boldsymbol{x}, \boldsymbol{y}) = \left[\sum_{i=1}^{\infty} (x_{i} - y_{i})^{2}\right]^{1/2}
\end{align*}
是 $X$ 上的一个度量。
\end{exercise}

\begin{exercise}\label{exercise:ch02-度量拓扑-11}
证明：若 $d$ 是 $X$ 的一个度量，则
\begin{align*}
d'(x, y) = d(x, y)/(1 + d(x, y))
\end{align*}
是诱导 $X$ 的拓扑的一个有界度量。[提示：对于 $x > 0$，设 $f(x) = x/(1 + x)$，应用中值定理证明 $f(a + b) - f(b) \leqslant f(a)$。]
\end{exercise}

\end{document}